\documentclass[12pt,a4paper]{article}

\usepackage[utf8]{inputenc}
\usepackage[T1]{fontenc}
\usepackage[english]{babel}
\usepackage{geometry}
\usepackage{graphicx}
\usepackage{hyperref}
\usepackage{enumitem}
\usepackage{listings}
\usepackage{xcolor}

\geometry{
    left=2.5cm,
    right=2.5cm,
    top=2.5cm,
    bottom=2.5cm
}

\hypersetup{
    colorlinks=true,
    linkcolor=black,
    urlcolor=blue
}

\lstset{
    basicstyle=\ttfamily\small,
    breaklines=true,
    frame=single,
    showstringspaces=false
}

\title{
    \textbf{Gym Management System}\\
    \large Database Management Systems Project
}

\author{
    Eda Bodo Ghislaine\\ 
    Matrikelnummer: 225786
}
\date{Sommersemester 2026}

\begin{document}

\maketitle
\tableofcontents
\newpage


% ============================================================
% 1. INTRODUCTION
% ============================================================

\section{Introduction}

The Gym Management System is a database application developed
for the Database Management Systems course.

The purpose of the application is to manage important information
related to a gym.

The system supports the management of:

\begin{itemize}
    \item members,
    \item courses,
    \item registrations,
    \item payments.
\end{itemize}

The application consists of a PostgreSQL database, a FastAPI
backend and a web-based frontend.


% ============================================================
% 2. SYSTEM ARCHITECTURE
% ============================================================

\section{System Architecture}

The application follows a three-layer architecture:

\begin{enumerate}
    \item Frontend
    \item Backend
    \item Database
\end{enumerate}

The frontend communicates with the backend using HTTP requests.

The backend processes the requests and accesses the PostgreSQL
database.

The general communication flow is:

\begin{center}
Frontend $\rightarrow$ FastAPI Backend $\rightarrow$ PostgreSQL
\end{center}

The backend and the PostgreSQL database are executed using
Docker Compose.


% ============================================================
% 3. TECHNOLOGIES% ============================================================
%===============================================================
\section{Technologies Used}

The following technologies are used in the project:

\begin{itemize}
    \item PostgreSQL
    \item FastAPI
    \item Python
    \item SQLAlchemy
    \item Docker
    \item Docker Compose
    \item HTML
    \item CSS
    \item JavaScript
    \item Git
    \item GitHub
    \item Debian packaging
\end{itemize}


% ============================================================
% 4. DATABASE
% ============================================================

\section{Database Design}

The application uses a PostgreSQL relational database.

The database contains four main tables:

\begin{itemize}
    \item members
    \item courses
    \item registrations
    \item payments
\end{itemize}


\subsection{Members}

The \texttt{members} table stores the members registered in
the gym.

The table contains:

\begin{itemize}
    \item \texttt{id}: primary key
    \item \texttt{name}: name of the member
    \item \texttt{email}: email address of the member
\end{itemize}

The email address is unique. Therefore, two members cannot
use the same email address.


\subsection{Courses}

The \texttt{courses} table stores the courses offered by the
gym.

The table contains:

\begin{itemize}
    \item \texttt{id}: primary key
    \item \texttt{name}: name of the course
    \item \texttt{trainer}: name of the trainer
    \item \texttt{schedule}: schedule of the course
\end{itemize}


\subsection{Registrations}

The \texttt{registrations} table connects members with
courses.

The table contains:

\begin{itemize}
    \item \texttt{id}: primary key
    \item \texttt{member\_id}: foreign key referencing
    \texttt{members.id}
    \item \texttt{course\_id}: foreign key referencing
    \texttt{courses.id}
\end{itemize}

The combination of \texttt{member\_id} and
\texttt{course\_id} is unique. This prevents the same member
from being registered more than once for the same course.


\subsection{Payments}

The \texttt{payments} table stores payment information for
gym members.

The table contains:

\begin{itemize}
    \item \texttt{id}: primary key
    \item \texttt{member\_id}: foreign key referencing
    \texttt{members.id}
    \item \texttt{amount}: payment amount
    \item \texttt{payment\_date}: date of the payment
\end{itemize}

Each payment is associated with an existing member through
the \texttt{member\_id} foreign key.


\subsection{Database Relationships}

The database relationships are:

\begin{itemize}
    \item A member can have multiple payments.
    \item A member can register for multiple courses.
    \item A course can contain multiple registered members.
    \item The \texttt{registrations} table connects members
    and courses.
\end{itemize}

The main relational structure can therefore be represented
as:

\begin{center}
Members $\longrightarrow$ Registrations $\longleftarrow$ Courses
\end{center}

\begin{center}
Members $\longrightarrow$ Payments
\end{center}
   


% ============================================================
% 5. BACKEND
% ============================================================

\section{Backend}

The backend is implemented with FastAPI.

The API provides CRUD operations for the main entities.

CRUD means:

\begin{itemize}
    \item Create
    \item Read
    \item Update
    \item Delete
\end{itemize}

The API contains routes for:

\begin{itemize}
    \item \texttt{/members/}
    \item \texttt{/courses/}
    \item \texttt{/registrations/}
    \item \texttt{/payments/}
\end{itemize}


% ============================================================
% 6. API SECURITY
% ============================================================

\section{API Security}

Write operations are protected using an API key.

The HTTP header used by the application is:

\begin{lstlisting}
X-API-Key
\end{lstlisting}

The operations POST, PUT and DELETE require the correct API key.

GET operations can be executed without an API key.

If the API key is missing or invalid, the server returns:

\begin{lstlisting}
401 Unauthorized
\end{lstlisting}


% ============================================================
% 7. FRONTEND
% ============================================================

\section{Frontend}

The frontend is implemented using HTML, CSS and JavaScript.

It provides sections for:

\begin{itemize}
    \item Members
    \item Courses
    \item Registrations
    \item Payments
\end{itemize}

The frontend communicates with the FastAPI backend using the
JavaScript \texttt{fetch()} function.

The backend runs on:

\begin{lstlisting}
http://localhost:8000
\end{lstlisting}

The frontend development server runs on:

\begin{lstlisting}
http://localhost:8080
\end{lstlisting}


% ============================================================
% 8. DOCKER
% ============================================================

\section{Docker Environment}

Docker Compose is used to start the backend and PostgreSQL
database.

The main services are:

\begin{itemize}
    \item \texttt{backend}
    \item \texttt{db}
\end{itemize}

The backend communicates with PostgreSQL through Docker's
internal network.


% ============================================================
% 9. DEBIAN PACKAGE
% ============================================================

\section{Debian Package}

The frontend can be packaged as a Debian package.

The generated package is called:

\begin{lstlisting}
gym-management-frontend.deb
\end{lstlisting}

After installation, the frontend files are stored in:

\begin{lstlisting}
/usr/share/gym-management/
\end{lstlisting}

The installed files are:

\begin{itemize}
    \item \texttt{index.html}
    \item \texttt{style.css}
    \item \texttt{app.js}
\end{itemize}


% ============================================================
% 10. TESTING
% ============================================================

\section{Testing}

The project contains automated API tests implemented in
\texttt{tests/test\_api.py}.

The tests verify the main API endpoints and the API key security
mechanism.

The following read operations are tested:

\begin{itemize}
    \item GET \texttt{/members/}
    \item GET \texttt{/courses/}
    \item GET \texttt{/registrations/}
    \item GET \texttt{/payments/}
\end{itemize}

All these endpoints must return:

\begin{lstlisting}
200 OK
\end{lstlisting}

The security tests verify that a POST request without an API key
or with an invalid API key returns:

\begin{lstlisting}
401 Unauthorized
\end{lstlisting}

A POST request with the correct API key must return:

\begin{lstlisting}
200 OK
\end{lstlisting}

The automated test then performs a complete CRUD test on a
temporary member:

\begin{enumerate}
    \item Create a temporary member using POST.
    \item Read the created member using GET.
    \item Update the member using PUT.
    \item Delete the member using DELETE.
    \item Verify that the deleted member can no longer be found.
\end{enumerate}

After deletion, requesting the same member must return:

\begin{lstlisting}
404 Not Found
\end{lstlisting}

This verifies both the CRUD functionality and the protection of
write operations using the \texttt{X-API-Key} header.


% ============================================================
% 11. VERSION CONTROL
% ============================================================

\section{Version Control}

Git is used for version control.

The project is stored in a GitHub repository.

The repository contains the following main directories:

\begin{itemize}
    \item \texttt{backend/}
    \item \texttt{frontend/}
    \item \texttt{packaging/}
    \item \texttt{tests/}
    \item \texttt{docs/}
\end{itemize}

Sensitive configuration information is stored in the
\texttt{.env} file.

The \texttt{.env} file is excluded from the Git repository
using \texttt{.gitignore}.


\subsection{Continuous Integration}

GitHub Actions is used to automatically build the project
documentation.

The workflow is stored in:

\begin{lstlisting}
.github/workflows/build.yml
\end{lstlisting}

For every push to the main branch, the workflow builds the
LaTeX documentation and verifies that the PDF was generated
successfully.

The documentation can also be attached automatically to a
GitHub Release when a version tag is created.

Version \texttt{v1.0} was created as the first release of the
project.


% ============================================================
% 12. CONCLUSION
% ============================================================

\section{Conclusion}

The Gym Management System demonstrates the interaction between
a relational database, a REST API and a web frontend.

The application allows the management of members, courses,
registrations and payments.

Docker simplifies the execution of the backend and database,
while the API key protects write operations.

The project also includes automated tests, Debian packaging and
Git version control.

\end{document}
